\section*{Models and border width (original Sections1 and3)}
An ABP is a layered directed graph with $d+1$ layers of vertices (edges only go from layers $i$ to $i+1$) such that the first and last layer have exactly the same number of vertices. Each vertex in the first layer has exactly one so-called corresponding vertex in the last layer. One interesting classical case is when the first and last layer have exactly one vertex, which is usually studied in theoretical computer science. We call this the single-(source,sink) model. Among algebraic geometers working on ABPs it is common to not impose restrictions on the number of vertices in the first and last layer. We call this the trace model. Every edge in an ABP is labeled by a homogeneous linear form. If we denote by $\ell(e)$ the homogeneous linear form of edge $e$, then we say that the ABP computes $\sum_p\prod_{e\in p}\ell(e)$, where the sum is over all paths that start in the first layer and end in the last layer at the vertex corresponding to the start vertex.
The width of an ABP is the number of vertices in its largest layer. We denote by $w(f)$ the minimal width required to compute $f$ in the trace model and we call $w(f)$ the trace ABP width complexity of $f$. We denote by $w_1(f)$ the minimal width required to compute $f$ in the single-(source,sink) model and we call $w_1(f)$ the single-(source,sink) ABP width complexity of $f$.
The border trace ABP width complexity $\underline w(f)$ is the smallest $w$ such that $f$ can be approximated arbitrarily closely by polynomials $f_\varepsilon$ with $w(f_\varepsilon)\le w$. Analogously, we define the border single-(source,sink) ABP width complexity $\underline w_1(f)$ as the smallest $w$ such that $f$ can be approximated arbitrarily closely by polynomials $f_\varepsilon$ with $w_1(f_\varepsilon)\le w$.
In the noncommutative homogeneous setting, $R_d=\mathbb F\langle x_1,\ldots,x_m\rangle_d$. We write $\ncw$ instead of $w$ and $\ncw_1$ instead of $w_1$ to highlight that we are in the noncommutative setting. The coefficient vectors identify this space with the $d$-th tensor power of $\mathbb F^m$. Over $\C$ the Zariski-closure of a set of polynomials of bounded complexity equals its Euclidean closure, by the standard principle cited in originalSection1.
For a homogeneous degree-$d$ ABP $\Gamma$, we denote by $V$ the set of vertices of $\Gamma$ and by $V^i$ the set of vertices in layer $i$, $1\le i\le d+1$. We choose an explicit bijection between the sets $V^1$ and $V^{d+1}$, so that each vertex $v$ in $V^1$ has exactly one corresponding vertex $\operatorname{corr}(v)$ in $V^{d+1}$. We denote by $E^i$ the set of edges from $V^i$ to $V^{i+1}$. Let $E$ denote the union of all $E^i$.
There is a classical interpretation in terms of iterated matrix multiplication: Fix some arbitrary ordering of the vertices within each layer, such that the $i$-th vertex in $V^1$ corresponds to the $i$-th vertex in $V^{d+1}$. For $1\le k\le d$ let $M_k$ be the $|V^k|\times|V^{k+1}|$ matrix whose entry at position $(i,j)$ in $M_k$ is the label from the $i$-th vertex in $V^k$ to the $j$-th vertex in $V^{k+1}$. Then $\Gamma$ computes the trace
\[\sum_{\substack{1\le k_1\le|V^1|\\1\le k_2\le|V^2|\\\cdots\\1\le k_d\le|V^d|}}(M_1)_{k_1,k_2}\cdots(M_{d-1})_{k_{d-1},k_d}(M_d)_{k_d,k_1}=\operatorname{tr}(M_1M_2\cdots M_d).\tag{8}\]
Hence the name trace model. In the single-(source,sink) model, the trace is taken of a $1\times1$ matrix.
\section*{5 Explicit construction of $f_0$ with higher complexity than border complexity}
Fix some $d\ge3$. In this section for every $m\ge2$ we construct $f_0$ such that
\[m=\underline{\ncw}(f_0)<\ncw(f_0).\tag{10}\]
A completely analogous construction can be used to find $f_0$ with $\underline w(f_0)<w(f_0)$ and with $\underline{\mathrm{mw}}(f_0)<\mathrm{mw}(f_0)$. For the sake of simplicity, we carry out only the proof for (10).
Recall that in a format $w$ ABP we have $w_{d+1}=w_1$. In each layer $i$ we enumerate the vertices $V^i=\{v^i_1,\ldots,v^i_{w_i}\}$ and we assume without loss of generality that the correspondence bijection between $V^{d+1}$ and $V^1$ is the identity on the indices $j$ of $v^1_j$, i.e., the $j$th vertex in $V^1$ corresponds to the $j$th vertex in $V^{d+1}$.
Fix an ABP format $w=(w_1,w_2,\ldots,w_d)$ such that for all $i$, $w_i\ge2$. Let $\Gcom$ denote the directed acyclic graph underlying an ABP of format $w$. An edge can be described by the triple $(a,b,i)$, where $1\le i\le d$, $1\le a\le w_i$ and $1\le b\le w_{i+1}$. Consider the following labeling of the edges with triple-indexed variables: $\ell_{\rm com}((a,b,i))=x^{(i)}_{(a,b)}$. Define $\fcom$ to be the polynomial computed by $\Gcom$ with edge labels $\ell_{\rm com}$.
We now construct $f_0$ as follows. Let $d$ be odd (the case when $d$ is even works analogously). Since in each layer we enumerated the vertices, we can now assign to each vertex its parity: even or odd. We call an edge between two even or two odd vertices parity preserving, while we call the other edges parity changing. Let us consider the following labeling of $\Gcom$: We set $\ell_0((a,b,i)):=x^{(i)}_{(a,b)}$ if $(a,b,i)$ is parity changing (i.e., $a\not\equiv b\pmod2$) and set the label $\ell_0((a,b,i)):=\varepsilon x^{(i)}_{(a,b)}$ otherwise, where $\varepsilon\in\C$. Let $f'_\varepsilon$ be the polynomial computed by $\Gcom$ with edge labels $\ell_0$ and set $f_\varepsilon:=\frac1\varepsilon f'_\varepsilon$ for $\varepsilon\ne0$. We define $f_0:=\lim_{\varepsilon\to0}f_\varepsilon$ (convergence follows from the construction, because $d$ is odd). By definition, for all $\varepsilon\ne0$, $f_\varepsilon$ can be computed by a format $w$ ABP. However, we will now prove that this property fails for the limit point $f_0$.
\setcounter{theorem}{3}
\begin{theorem}\label{thm:nonclosure}
Fix an ABP format $w=(w_1,w_2,\ldots,w_d)$ such that for all $i$, $w_i\ge2$. Let $f_0$ be defined as above. Then, $f_0$ cannot be computed by an ABP of format $w$.
\end{theorem}
Note that for a format where $m=w_1=\cdots=w_d$, this gives the $f_0$ which was desired in (10). (Note, however, that $f_0$ can be computed by an ABP of width $2m$ as follows. Construct an ABP $\Gamma'$ that has, for each vertex $v\in\Gcom$, vertices $v'$ and $v''$. For each parity changing edge $(a,b)\in\Gcom$ with label $\ell_0$, add edges $(a',b')$ and $(a'',b'')$ with the same label $\ell_0$. For each parity preserving edge $(a,b)\in\Gcom$ with label $\ell_0$, add edge $(a',b'')$ with label $(\frac1\varepsilon)\ell_0$. For corresponding vertices $u,v$ in $\Gcom$, let $v''$ be the corresponding vertex for $u'$ and $v'$ be the corresponding vertex for $u''$ in $\Gamma'$. All paths between corresponding vertices in this ABP use exactly one parity preserving edge of $\Gcom$, and so this ABP computes $f_0$.)
The proof of Theorem 4 works as follows. Let $G:=\GL_{w_1w_2}\times\GL_{w_2w_3}\times\cdots\times\GL_{w_dw_{d+1}}$. Let $\End:=\overline G$ denote its Euclidean closure, i.e., tuples of matrices in which one or several matrices can be singular.
We consider noncommutative homogeneous polynomials in the variables $x^{(i)}_{(a,b)}$ such that the $i$-th variable in each monomial is $x^{(i)}_{(a,b)}$ for some $a\in[w_i]$ and $b\in[w_{i+1}]$. The vector space of these polynomials is isomorphic to $W:=\C^{w_1w_2}\otimes\C^{w_2w_3}\otimes\cdots\otimes\C^{w_dw_{d+1}}$ and the monoid $\End$ (and thus also the group $G$) acts on this space in the canonical way. The set
\[\{f\in W\mid f\text{ can be computed by a format }w\text{ ABP}\}\]
is precisely the orbit $\End\fcom$. We follow the overall proof strategy in [24]. The monoid orbit $\End\fcom$ decomposes into two disjoint orbits:
\[\End\fcom=G\fcom\cup(\End\setminus G)\fcom.\]
Our goal is to show two things independently:
\begin{enumerate}\item $f_0\notin(\End\setminus G)\fcom$, and\item $f_0\notin G\fcom$,\end{enumerate}
which finishes the proof of Theorem 4.
All elements in $(\End\setminus G)\fcom$ are not concise, a term that we define in Section 6, where we also prove that $f_0$ is concise. Therefore $f_0\notin(\End\setminus G)\fcom$.
All elements in $G\fcom$ have full orbit dimension, a term that we define in Section 7 and we prove that $f_0$ does not have full orbit dimension in Section 8. This finishes the proof of Theorem 4.
